% Generated draft: requires mathematical and source review.
\documentclass[11pt]{article}
\usepackage[margin=25mm]{geometry}
\usepackage{fontspec}
\setmainfont{Latin Modern Roman}
\usepackage{amsmath,amssymb,mathtools,unicode-math}
\setmathfont{Latin Modern Math}
\usepackage{enumitem,needspace,xurl,hyperref}
\setlength{\parindent}{0pt}
\setlength{\parskip}{5pt}
\setlength{\emergencystretch}{3em}
\newcommand{\N}{\mathbb{N}}
\newcommand{\Z}{\mathbb{Z}}
\newcommand{\Q}{\mathbb{Q}}
\newcommand{\R}{\mathbb{R}}
\newcommand{\C}{\mathbb{C}}
\newcommand{\F}{\mathbb{F}}
\newcommand{\A}{\mathbb{A}}
\newcommand{\PP}{\mathbb{P}}
\newcommand{\E}{\mathbb{E}}
\newcommand{\eps}{\varepsilon}
\newcommand{\dd}{\,\mathrm{d}}
\begin{document}
\section*{UnsolvedMath ID: 30000154}
\textbf{Source problem code:} catalogue entry
\subsection{Definitions and mathematical statement}
For a prime $p$, a set $A\subset\Z/p\Z$ is sum-free if $(A+A)\cap A=\varnothing$, allowing the two summands to coincide. Its density is $|A|/p$. For an integer $1\le n\le p/2$, let
\[
 I_{p,n}=\{n,n+1,\ldots,p-n\}\pmod p.
\]
A dilation is multiplication by $u\in(\Z/p\Z)^\times$; no translation is included.

The source asks for the sharp structural threshold, between the historical cardinality scales $0.2p$ and $0.33p$, above which every sum-free set is contained in some $uI_{p,n}$. The dimensionless densities are $0.2$ and $0.33$; the factors $p$ refer to set sizes. The decimal $0.33$ is retained and not silently replaced by $1/3$.

The extract does not define $n$. A natural intended strong version is $n=|A|$, leading to the property
\[
 |A|>\gamma p\quad\Longrightarrow\quad
 \exists u\in(\Z/p\Z)^\times,\ A\subset uI_{p,|A|}
\]
for all sufficiently large primes and all sum-free $A$. In that asymptotic interpretation, one seeks the infimum of constants $\gamma$ for which this property holds uniformly. This is an explicit interpretation, not a verified recovery of the omitted definition of $n$ or of strict versus weak threshold endpoints. If $n$ were freely chosen, $n=1$ would make containment essentially trivial since a nonempty sum-free set excludes zero; thus its relation to $A$ is essential. The source also does not settle exact finite-prime versus asymptotic quantifiers.

The literature assessment checked on 21 August 2026 says the original $0.33$ threshold was improved in later work and mentions lower examples near $0.2$, while the full structural threshold remains unresolved in that assessment. No improved numerical constant or exact lower bound is supplied. The question concerns the stated interval-containment structure, not merely the maximal possible size of a sum-free set. These are dated triage claims requiring independent verification.
\subsection{Short English statement}
Find the sharp density forcing a sum-free set modulo a prime into the specified dilated central interval. The interval parameter must be tied to the set, presumably its size; otherwise the containment condition loses its intended structural force.
\subsection{Sources}
\begin{itemize}
\item \emph{Finite Fields: Theory and Applications} (2005), Oberwolfach Reports, DOI 10.4171/owr/2004/54, pp. 2913--2970, \url{https://doi.org/10.4171/owr/2004/54}; proposed year supplied as 2004.
\item Literature assessment checked 21 August 2026: original OWR report, \url{https://doi.org/10.4171/owr/2004/54}; Candela--Gonzalez-Sanchez--Grynkiewicz, \url{https://doi.org/10.24033/bsmf.2827}.
\item Local UnsolvedMath registry, permanent ID $30000154$, exact source code \texttt{catalogue entry}, created 25 August 2026 and updated 16 September 2026. Corrected-and-verified status records balanced group, density and interval notation; it does not supply the missing relation defining $n$.
\end{itemize}
\end{document}
